% FORMALIZACION_REUNIDA: CG-015/P09 del inventario REV05.
% Propietarios: sections/nucleo.tex (Cayley APP), libro/toro27.tex:123--256,
% PAQUETE_ARTICULO_I_INTEGRACION_ESTADO_CAMPOS_20260919/deltas/integracion/
% APPRouteWinding.lean (desplazamiento y cociclo entero).
% Esta insercion explicita el cambio de soporte y las fibras aritmeticas.
\section{Raffinement bidimensionnel de l'APP et cocycle de retenue}
\label{sec:refinamiento-app-27-9}

La réalisation de l'APP sur un tore de neuf positions par coordonnée
et son extension à vingt-sept positions sont reliées par la division
euclidienne. L'extension conserve deux directions indépendantes et attribue
un chiffre ternaire supplémentaire à chaque coordonnée. Ainsi, chaque position
de l'APP initiale possède neuf relèvements dans le support étendu.
Les deux chiffres supplémentaires conservent la position dans la fibre,
tandis que la direction du transport et le régime du lecteur TRIT
demeurent des coordonnées propres de l'état enrichi.

Pour $m\geq1$, nous écrivons
\[
 T_m=(\mathbb Z/m\mathbb Z)^2.
\]
Dans les calculs de cette section, on choisit les représentants
$0,\ldots,m-1$. La présentation nonadique $1,\ldots,9$ imprime la
classe zéro au moyen du signe $9$. Cette convention permet de lire la table
APP et d'effectuer les divisions entières avec une règle unique de résidus.

\subsection{Projection, fibres et coordonnées ternaires de raffinement}

La projection et la section de représentants sont
\[
 p:T_{27}\longrightarrow T_9,\qquad
 p(x,y)=(x\bmod9,y\bmod9),
\]
\[
 s:T_9\longrightarrow T_{27},\qquad s(r,t)=(r,t),
 \quad 0\leq r,t<9.
\]
La projection est un homomorphisme surjectif de groupes additifs. Son
noyau est
\[
 \ker p=\{(9a,9b):a,b\in\mathbb F_3\}
        \simeq\mathbb F_3^2.
\]
L'application de coordonnées
\begin{equation}
 \Psi:T_9\times\mathbb F_3^2\longrightarrow T_{27},\qquad
 \Psi((r,t),(a,b))=(r+9a,t+9b)
 \label{eq:app-refinamiento-coordenadas}
\end{equation}
utilise les représentants $0\leq a,b<3$.

\begin{proposition}[Décomposition unique des positions étendues]
\label{prop:app-refinamiento-fibras}
L'application $\Psi$ est bijective. Chaque fibre de $p$ contient neuf
positions et la totalité du support satisfait
\[
 |T_{27}|=27^2=9^2\cdot3^2=729.
\]
Les deux coordonnées de $\mathbb F_3^2$ sont les deux trits positionnels
qui décrivent le raffinement des deux coordonnées de l'APP.
\end{proposition}
\begin{proof}
Pour $0\leq x,y<27$, les divisions euclidiennes
\[
 x=r+9a,\qquad y=t+9b,
 \quad 0\leq r,t<9,\quad0\leq a,b<3
\]
existent et sont uniques. Ces formules construisent l'inverse de $\Psi$.
Une fois $r,t$ fixés, les neuf paires $(a,b)$ sont distinctes et épuisent la
fibre. Le recensement résulte de la multiplication des $81$ positions résiduelles
par leurs neuf relèvements.
\end{proof}

Par exemple, la position $(5,5)$ possède la fibre
\[
 \{(5+9a,5+9b):0\leq a,b<3\}.
\]
La ligne de cette fibre avec $b=1$ est
$(5,14),(14,14),(23,14)$. Chaque position conserve la même paire résiduelle
$(5,5)$ et une paire distincte de quotients. La multiplicité neuf exprime
deux choix ternaires indépendants. La lecture topologique et
temporelle du TRIT conserve, en outre, son propre domaine opératoire.

\subsection{Composition additive et mémoire du franchissement modulaire}

Pour $r,t\in\{0,\ldots,8\}$, nous définissons
\[
 r\oplus t=(r+t)\bmod9,\qquad
 c_9(r,t)=\left\lfloor\frac{r+t}{9}\right\rfloor.
\]
Transportée par $\Psi$, l'addition dans une coordonnée de
$\mathbb Z/27\mathbb Z$ est
\begin{equation}
 (r,a)\boxplus(t,b)
   =\bigl(r\oplus t,\ a+b+c_9(r,t)\pmod3\bigr).
 \label{eq:app-refinamiento-adicion}
\end{equation}
En deux coordonnées, la même formule s'applique séparément à chaque axe.

\begin{proposition}[Cocycle de l'extension positionnelle]
\label{prop:app-refinamiento-cociclo}
La retenue satisfait
\begin{equation}
 c_9(r,t)+c_9(r\oplus t,u)
 =c_9(t,u)+c_9(r,t\oplus u)
 =\left\lfloor\frac{r+t+u}{9}\right\rfloor.
 \label{eq:app-refinamiento-cociclo}
\end{equation}
Par conséquent, \eqref{eq:app-refinamiento-adicion} est associative et
reproduit exactement le groupe additif du support étendu.
\end{proposition}
\begin{proof}
L'égalité $r+t=(r\oplus t)+9c_9(r,t)$ donne
\[
 \left\lfloor\frac{r+t+u}{9}\right\rfloor
 =c_9(r,t)+\left\lfloor\frac{(r\oplus t)+u}{9}\right\rfloor.
\]
Le regroupement $r+(t+u)$ fournit l'autre membre. Substituer
$x=r+9a$ et $y=t+9b$ dans $x+y$ prouve
\eqref{eq:app-refinamiento-adicion} ; la réduction de sa seconde
composante modulo trois correspond à réduire $x+y$ modulo vingt-sept.
L'identité du cocycle prouve directement l'associativité de cette
écriture par résidus et quotients.
\end{proof}

En particulier, $8+1=0+9\cdot1$ montre le franchissement élémentaire : le résidu
revient à la classe zéro, imprimée comme $9$ dans la table nonadique, et le
premier quotient augmente d'une unité. Si l'on conserve uniquement le
résidu, on perd la donnée qui distingue $0$, $9$ et $18$ dans le support
étendu ; l'état relevé distingue ces trois positions.

L'extension
\[
 0\longrightarrow\mathbb Z/3\mathbb Z
 \xrightarrow{\ a\mapsto9a\ }
 \mathbb Z/27\mathbb Z
 \longrightarrow\mathbb Z/9\mathbb Z\longrightarrow0
\]
est non scindée en tant qu'extension de groupes. En effet, tout
relèvement de la classe $1$ est $1+9a$, d'ordre vingt-sept. Une
section homomorphe devrait envoyer un élément d'ordre neuf sur un élément
dont l'ordre divise neuf. La section de représentants possède, en
revanche, le défaut exact $9c_9$, que la formule précédente conserve.

\subsection{Compatibilité du feuillet multiplicatif}

La même décomposition conserve la seconde opération de l'APP. Pour
une coordonnée, nous définissons
\[
 r\odot t=rt\bmod9,\qquad
 m_9(r,t)=\left\lfloor\frac{rt}{9}\right\rfloor.
\]
Le produit relevé est
\begin{equation}
 (r,a)\boxtimes(t,b)
 =\bigl(r\odot t,\ m_9(r,t)+rb+ta\pmod3\bigr).
 \label{eq:app-refinamiento-producto}
\end{equation}
\begin{proof}
Le développement entier fournit
\[
 (r+9a)(t+9b)=rt+9(rb+ta)+81ab.
\]
Le dernier terme s'annule modulo vingt-sept. En substituant
$rt=(r\odot t)+9m_9(r,t)$, on obtient la formule. L'unicité de la
division euclidienne prouve que ses deux composantes récupèrent le produit
résiduel et son quotient de raffinement.
\end{proof}
Par exemple, $8\cdot8=64\equiv10\pmod{27}$ est représenté par
$(1,1)$ : le résidu est $1$ et
$m_9(8,8)=7\equiv1\pmod3$. Ainsi, somme et produit conservent leurs
retenues respectives au sein d'un même support étendu.

\subsection{Directions et réalisation du transport}

Soit $\mathcal D=\{N,S,E,O\}$ le registre de directions et
$\delta_d\in\mathbb Z^2$ le déplacement unitaire associé à $d$.
Sur le support avec direction, nous définissons
\[
 S_m(z,d)=(z+\delta_d\pmod m,d).
\]
La projection conserve explicitement la direction,
$p_{\mathcal D}(z,d)=(p(z),d)$, et satisfait
\begin{equation}
 p_{\mathcal D}S_{27}=S_9p_{\mathcal D}.
 \label{eq:app-refinamiento-direccion}
\end{equation}
Le calcul du quotient utilise, pour un déplacement entier $v$,
\[
 c_9(r;v)=\left\lfloor\frac{r+v}{9}\right\rfloor,
 \qquad r'=r+v-9c_9(r;v),\qquad
 a'=a+c_9(r;v)\pmod3.
\]
La formule inclut les déplacements négatifs. Par exemple,
$r=0,v=-1$ donne $r'=8,c_9=-1$ ; le retour inverse restitue le quotient.

\begin{proposition}[Conjugaison du déplacement en coordonnées de fibre]
Le changement de coordonnées $\Psi$ transforme $S_{27}$ en le
déplacement résiduel $S_9$ accompagné de la mise à jour des
deux quotients qui vient d'être définie. Cette transformation est bijective
et conserve la direction. Ses puissances conservent la même identité
pour toute longueur finie du transport.
\end{proposition}
\begin{proof}
Sur chaque axe, on substitue $x=r+9a$ dans $x+v$ et l'on effectue la division
euclidienne. La direction demeure fixe dans la définition de $S_m$.
L'inverse utilise $-v$ et le cocycle, ce qui permet de récupérer l'état
initial. Composer l'identité une fois par transition prouve
l'énoncé pour toutes les puissances.
\end{proof}

La réalisation de Hilbert conserve cette décomposition. Sur les bases
de positions, l'opérateur
\[
 \mathcal U\delta_{\Psi(z,a),d}
       =\delta_{z,d}\otimes\delta_a
\]
définit une isométrie surjective
\[
 \ell^2(T_{27}\times\mathcal D)
 \simeq\ell^2(T_9\times\mathcal D)\otimes\ell^2(\mathbb F_3^2).
\]
Le déplacement conjugué est l'opérateur de permutation avec retenue
décrit dans la proposition. Un facteur interne supplémentaire $E$ demeure
intact lors du produit tensoriel avec $I_E$. Si un lecteur de phase dépend seulement de
la projection et qu'un opérateur directionnel agit avec la même matrice à
toutes les positions, leurs carrés de compatibilité
commutent également. Un lecteur qui utilise le quotient se réalise, précisément,
dans le facteur de fibre conservé. Ces conditions explicitent la
relation entre le support étendu et les facteurs de l'opérateur sur
le tore, avec les actions internes déclarées dans sa construction.

\subsection{Six chiffres ternaires et structures algébriques distinctes}

Chaque coordonnée de $T_{27}$ possède trois chiffres en base trois. Leur écriture
conjointe définit une bijection d'ensembles
\[
 T_{27}\longleftrightarrow\mathbb F_3^6.
\]
Cette bijection publie six chiffres positionnels. L'addition du tore
contient des retenues entre les chiffres et l'addition vectorielle de
$\mathbb F_3^6$ s'effectue composante par composante. Par exemple, la
somme positionnelle $2+1=3$ transforme $(2,0,0)$ en $(0,1,0)$, tandis
que la somme vectorielle ternaire donnerait $(0,0,0)$. Le tore contient
des éléments d'ordre vingt-sept ; l'espace vectoriel ternaire a pour
exposant trois. Les deux objets partagent un codage fini,
avec des opérations typées différentes. La fibre de mots ternaires,
le support de positions et les droites d'une surface cubique
conservent ainsi leurs constructions et leurs applications spécifiques.
